% RENUMBERED 2026-11-16 (user: "rename the files also"): methodology files follow the reading order.
%   Old -> new: 05 features -> 03 (05a/05b/05c -> 03a/03b/03c), 06 feature validity -> 04,
%   03 statistical model (before that 03_prediction_model) -> 05, 04 model size -> 06; 07 unchanged.
%   File names in the comments below were updated; bare numbers in comments dated before this entry
%   ("Methods 03", "05a header", "05d") are the OLD numbers.
% Methods, subsection of sec:method-features: the two diameter-ratio features only. Why the bifurcation
% is represented at all belongs to sec:method-features and sec:lit-wss; which segment rows they enter
% (proximal LAD and LCx) to 03_features.tex and sec:method-models.
% 2026-11-11: written from daugther_ration_plan.md (same folder). Replaces the notes trimmed 2026-11-07
% (full draft in git, commit e4eb857).
% Decision (with the user, 2026-11-11): two features, still 2 parameters: daughter asymmetry
%   ln(D_LCx / D_LAD) and Finet's ratio D_LM / (D_LAD + D_LCx). Angle and the two ratios keep their own
%   linear terms in the tested geometry block (5 parameters); angle x asymmetry is one prespecified
%   secondary term, estimated, outside the LR test (one simulation study on synthetic trees, and a
%   product needs its main effects). Owned by sec:method-models.
% Rejected, not stated in the text (counterfactual rule): the LM->LAD and LM->LCx ratios planned
%   before 2026-11-11. Both divide by the LM diameter, the least reliable of the three (LM length
%   10.5 +- 5.3 mm in Medrano-Gracia 2016, so the stretch clear of the ostium and of the confluence is
%   short or absent); its error enters both and correlates them, and both go missing together. The
%   review supports the LAD-LCx width difference (Garcha 2025) and discusses Finet's ratio (Zhang 2024,
%   null); neither used LM->daughter ratios.
% Murray's exponent 3, the pooled 2.39 (taylor2024murray) and 7/3 (huo2012scaling) are no longer
%   needed here; the fact-ownership line for them in the content map (Thesis_coronary/CLAUDE.md) lapses.
% Implementation: CGPS_data/thesis_final/daughter_ratio.py (contour_area, contour_diameter,
%   reference_diameter, lm_ratios, tree_ratios; DEPTHS = (5, 10) mm, the interval of
%   bifurcation_angle.DEPTHS). Checked against the text line by line 2026-11-11: window 5 to 10 mm
%   from O, upstream for the LM; median over the points in the window; nan when the vessel does not
%   reach 10 mm or fewer than 2 points have a diameter there; no shared O -> both nan; short LM ->
%   Finet nan, asymmetry kept. Not in the text: the 3D polygon area (cross products about the centroid).
% Figure: content/methodology/figures/surface_cut_figure.py since 2026-11-19 (user: replace fig:method-ratio;
%   made-up lumen meshes run through daughter_ratio.section, contour_area, contour_diameter). The old
%   contour-based daughter_ratio_figure.py is no longer included. CAREFUL: the figure shows plane cuts of
%   the lumen surface (the code since 2026-11-19), while the body text above still says the tracer gives a
%   contour at each centerline point and states the two-contour missing rule; the body must be rewritten
%   to the surface cuts and the strict 11-cut rule before this section is approved.
% 2026-11-19 (user: no panel letters, as fig:method-curvature and fig:method-angle): letters removed from
%   surface_cut_figure.py, the panel references (B and C, D) cut from the text, the caption names the steps
%   in order and now names p_i, p_{i+1}.
% 2026-11-19 (user: rewrite the body to the surface cuts). Matches daughter_ratio.tree_asymmetry,
%   daughter_asymmetry, surface_diameters and section (read 2026-11-19): the tracer lumen is a triangle mesh
%   (lumen_probe.md), cut by the plane through the centerline point perpendicular to the centerline at the 11
%   depths SECTION_DEPTHS (5 to 10 mm from O, every 0.5 mm); the closed loop enclosing the point is the
%   cross-section (the smallest if several); no diameter when no loop encloses the point or the loop also
%   encloses the other daughter's crossing; strict rule: D = median of the 11 only when all 11 have a diameter;
%   missing as the angle (no LM label, LAD or LCx missing, no shared O, a daughter under 10 mm). Resolves the
%   CHECK on the tracer contour format and the CAREFUL of the figure entry above.
%   Cut: "To keep closely spaced contours from counting more, D was read every 0.5 mm" (the cuts are at fixed
%   depths, so no interpolation); the two-contour missing rule (superseded by the strict rule).
%   Not in the text (code level, methods-not-code-level rule): the tangent from the points 1 mm before and
%   after, CUT_RADIUS 10 mm, the smallest of several enclosing loops, the triangle-fan area (in the figure).
%   Gharleghi still "adapted": their median is over 0 to 10 mm sampled every 0.01 mm (arXiv 2311.18489, 2.3,
%   reread 2026-11-19); how their cross-sections were cut is not stated there.
% CAREFUL: tree_asymmetry is not yet called by progression_who_pooled.py, which still builds the old
%   narrower-row coding; the model and this text disagree until it does.
% Sources: medranogracia2016atlas ("effective radius", A = pi R^2, read 2026-11-10 for 05b);
%   finet2008fractal (0.678 as the reported mean of D_mother / (D_d1 + D_d2)).
% CHECK: Garcha full text not read (Nature redirect); abstract: "diameter ratio between left coronary
%   branches ... can amplify the effect of ... bifurcation angle". Its exact ratio definition and
%   direction are unverified, so the secondary term is not described as a replication.
% CHECK: 0.678 taken from the plan, not reread in Finet et al. this session. The plan's arithmetic
%   "3.39 / (2.708 + 2.236)" gives 0.686, not 0.678 (a ratio of means is not a mean of ratios); the
%   code check uses tubes of 3.39, 2.75 and 2.25 mm instead.
% CHECK: tracer contour format unknown (daughter_ratio.tracer_rings raises), so "the tracer's lumen
%   contour at each centerline point" is an assumption until the files are seen.
% (2026-11-13: the earlier CHECK on a 5 to 10 mm LM window is void; the LM window is the first 2 mm
%   from the ostium, see the 2026-11-11 entries below. Current definition: D = median area-equivalent
%   diameter, LAD/LCx 5 to 10 mm from O, LM 0 to 2 mm from the ostium; a = ln(D_LCx/D_LAD),
%   F = D_LM/(D_LAD + D_LCx). Report the coverage of F once the tracer LM lengths are known.)
% CAREFUL: progression_who_pooled.py has no ratios yet, so the model and this text disagree until it
%   calls tree_ratios. (2026-11-18: its design() now has an "asymmetry" term built with
%   narrower_rows, but nothing fills the person column until tracer_rings reads tracer contours.)
% Style: prose only, British spelling, reason before method, no "trunk", no sentence opening "Because".
% 2026-11-11 (user: no First/Subsequently/Finally for short methods, true to the code but not every
%   detail, flow between sentences, A++ thesis voice). Body rewritten as connected prose; facts, labels,
%   cites and equations unchanged. Each step now follows from the one before (contour -> diameter ->
%   one diameter per vessel -> the two ratios -> missing rules). The stretch and median reasons merged
%   into one sentence after the method; "The cost is that a left main shorter than 10 mm has no
%   diameter" folded into the missing sentence. "A vessel with fewer than two contours counted as too
%   short" now stated as its effect (both ratios missing for a daughter, Finet only for the LM), as in
%   lm_ratios. Caption unchanged.
% 2026-11-11 (user: concise, only what fits a methodology section, as for 05b). Cut: "so D is the
%   diameter of the circle with the same area" ("area-equivalent" says it), the reading of a (0 for equal
%   daughters, negative when the LCx is narrower) and of F ("compares the left main with its two daughters
%   together", "a lower value means daughters that are wide for their parent"), and Finet's reported mean
%   0.678 (reference value, belongs in Results/Discussion when CGPS values are compared). Both features
%   now in one sentence with Figure B; "the same stretch as for the bifurcation angle" moved from the
%   opener to the stretch sentence; the LM-uncertainty reason joined into one sentence.
% 2026-11-11 (user: add a citation for the method): medranogracia2017bifurcation added as prior use of
%   Finet's ratio on CTA (read 2026-11-11, PMC5323506: D_PMV / (D_DMV + D_SB), 446 bifurcations in 211
%   normal CTAs, mean 0.6576 +- 0.083; the value is kept for Results/Discussion, as 0.678). CAREFUL: its
%   Methods do not say how the diameters were measured (lumen meshes cut by a 10 mm sphere at the
%   bifurcation point only), so cite it for the ratio on CTA, never for the diameter or the window. No
%   source found for the 5 to 10 mm diameter window or the median; both stay the thesis's own choices.
% 2026-11-11 (user: where is the LM diameter measured?): the stretch sentence split in two, LAD/LCx
%   downstream of O, LM upstream between O and the ostium. Checked medranogracia2016atlas (Computation
%   of Diameters, read 2026-11-11): effective diameter "every 0.5 mm from the ostium or previous main
%   bifurcation", i.e. the LM profile runs downstream from the ostium (Fig. 6A "Left main from ostia
%   (continuing to LAD)"), daughters from the bifurcation; LM reported at the ostium (3.7 mm, Results;
%   abstract 3.5 +- 0.8 mm, unexplained). No interval averaging for diameters (5 to 10 mm is for angles
%   only), and the positions behind their Finet's ratio are not stated. So the LM window upstream of O is
%   this thesis's choice, not theirs; never write "following" for it.
% 2026-11-11 (user: measure the LM at the ostium like the atlas; median over the first 2 mm, stated as a
%   design choice extending the atlas). Code: daughter_ratio.LM_OSTIUM = (0, 2) mm from the ostium, the
%   LM path no longer reversed; check() and the synthetic split_trunks tree confirm the window sits at the
%   ostium end. Missing rule now LM < 2 mm or < 2 contours there. Supersedes the 5 to 10 mm upstream
%   window and its CHECK above. The atlas Results give 3.7 mm at the ostium (abstract 3.5 +- 0.8 mm).
% CHECK: with an LM under about 4 mm the 2 mm window touches the confluence at O; and the number of
%   contours in 2 mm depends on the tracer point spacing (unknown). Report both once the tracer runs.
% CHECK: "where the lumen leaves the aorta" assumes the tracer's first LM contour sits at the aortic
%   junction; confirm on the tracer output.
% 2026-11-11 (user: add and cite Gharleghi 2025): gharleghi2025sex (IEEE TBME 72(7), 2025; read as
%   arXiv 2311.18489, 2.3): "The median diameter for each vessel, corresponding to the diameter of a
%   circle with equivalent cross-sectional area, was calculated over the 10mm segment" of each branch
%   within 10 mm arc length of the bifurcation point; Finet 0.57 in both sexes. "Adapted", not
%   "following": their stretch is 0 to 10 mm from the bifurcation point for all three vessels, ours 5 to
%   10 mm for the daughters and 0 to 2 mm from the ostium for the LM. Their reason for the median
%   (diameter changes little along 10 mm, SD 0.1 mm) differs from ours (side-branch origins).
% CHECK: bib entry from the arXiv author list and the journal citation found by search; pages and the
%   published author list not checked against IEEE Xplore.
% 2026-11-13 (user): Finet's ratio DROPPED as a model feature. Reasons (counterfactual rule, not in the
%   text): the review's only evidence is Zhang 2024's null; no sourced reason for a link to plaque
%   (Schoenenberger 2012, Murray deviation vs calcification, abstract only, cross-sectional); it is the
%   only feature needing the LM diameter, the least certain of the three, and it goes missing for a short
%   LM. Geometry block now 4 parameters (+1 secondary angle x asymmetry). The LM window, eq:method-finet,
%   finet2008fractal and medranogracia2017bifurcation are out of the text; the notes above on them are
%   history. daughter_ratio.py still returns "finet" (descriptive only, not in the model).
% 2026-11-18 (user): the angle x asymmetry term DROPPED (Garcha p. 7: the angle's effect changes sign
%   with which daughter is wider, so one product term does not test it); its sentence is cut. Geometry
%   block 3 terms, no secondary product.
% 2026-11-18 (user, plan approved): the asymmetry enters ONLY THE NARROWER DAUGHTER's proximal row as
%   |a|, the wider one gets 0 (pLAD = max(a, 0), pCX = max(-a, 0); daughter_ratio.narrower_rows, checked
%   in check()). eq:method-asymmetry keeps the signed a. Reason: Garcha p. 6 "ALWSS formation in the
%   smaller-of-the-two branches", Fig. 4 caption (stated in sec:lit-wss, pointed back to, not cited
%   again). Decided before the tracer cohort was linked to the outcome. Scaling (by SD, not centred)
%   is stated in sec:method-model-spec.
%   CAREFUL: synthetic rigid CFD trees, outcome low-WSS area, not plaque.
%   CHECK before outcome linkage (aggregate only): share of participants whose LCx is the narrower; if
%   it is most of them, the term mostly acts in the proximal LCx, which has fewer events. Atlas means
%   in Garcha (LAD 3.7 vs LCx 3.6 mm men, 3.0 vs 2.8 mm women, SD about 0.6) suggest both occur.
%   Not modelled: one slope for both daughters, though Garcha's low-WSS area was larger with a wider
%   LCx (16.7% vs 6.8%, p. 7). Near-equal daughters give |a| near 0, so which one counts as narrower
%   hardly matters.
% 2026-11-18 (code review, user: fix the issues): the median was over the contour POINTS in the stretch,
%   so densely sampled parts counted more. daughter_ratio.reference_diameter now reads D every 0.5 mm
%   (STEP, the depths of bifurcation_angle.DEPTHS) by linear interpolation between contours and takes the
%   median of those values; rule kept: nan below 2 contours in the stretch or a vessel short of 10 mm.
%   check(): crowded points no longer pull the median. One sentence added after the stretch sentence.
%   The left main path for Finet's ratio now uses split_trunks.shared_prefix, the rule split() uses for O.
%   Missing asymmetry: complete cases (user, 2026-11-18), stated in sec:method-model-spec.
% 2026-11-18 (user: two models; the angle and asymmetry go into a person-level cloglog GLM whose outcome is
%   plaque in the proximal LAD or proximal LCx). The narrower-row coding is superseded: |a| enters the
%   bifurcation model once per participant (per SD, not centred, sec:method-model-spec). The Garcha reason
%   (low WSS in the narrower daughter) is kept for taking |a|. CAREFUL: the outcome covers both daughters,
%   so the model no longer ties |a| to the narrower daughter's own segment; do not write that it does.
%   CODE: daughter_ratio.narrower_rows and the asymmetry rows in progression_who_pooled.design() still use
%   the old coding; no person-level GLM exists yet.
% 2026-11-20 (user: write the reason for the ln into the section): one sentence after eq:method-asymmetry. A plain
%   ratio puts "twice as wide" at 2 and 0.5 (distance 1 vs 0.5 from equal); the log puts both at ln 2 from 0,
%   so |a| treats a wider LAD and a wider LCx alike and the order of the ratio does not matter. Own reasoning,
%   no cite. Garcha (papers/s41598-025-85781-x (1).pdf, read 2026-11-20) uses R_D = D_LAD/D_LCx, set in
%   synthetic models, with "skewed dominance" at R_D <= 0.65 or >= 1.40 and a quadratic fit of ALWSS on R_D;
%   the paper never says "asymmetry" or "daughter", so the term stays the thesis's own (no "referred to as").
% 2026-11-20 (user, /plan approved: rewrite to the tone and depth of sec:method-angle). Opener names the feature and
%   points to the angle's eleven points instead of restating the stretch; steps in order (cut, diameter, median,
%   "Finally" the asymmetry); inserted clauses split; "referred to as a triangle mesh" dropped (term unused, the
%   caption explains the area). Facts, labels, equations and cites unchanged; checked against daughter_ratio.py
%   (SECTION_DEPTHS, contour_diameter, daughter_asymmetry strict rule, tree_asymmetry = the angle's rules).
%   medranogracia2016atlas reopened (~/Downloads/web139_EIJ-D-16-00111_MedranoGraciaCR104.pdf, Computation: "From
%   the area (A), effective radius (R) and diameter were determined (A = pi R^2)"): supports the diameter.
% CHECK: gharleghi2025sex (median over 0 to 10 mm, "adapted") could not be reopened on 2026-11-20: no local PDF;
%   put it in papers/ to confirm before submission.
% 2026-11-20 (user: inform the reader first how the diameter is estimated). Order now: opener, diameter (cut,
%   area-equivalent D, the eleven points, median), then the asymmetry. Medrano-Gracia et al. credited for the whole
%   diameter estimate: Methods, "Computation of diameters": "the mesh was intersected with a plane normal to the
%   centreline ... polygon shape from which the area was calculated. From the area (A), effective radius (R) and
%   diameter were determined (A = pi R^2)", every 0.5 mm (reopened 2026-11-20). Not theirs: the 5 to 10 mm stretch for
%   diameters (they use it for angles), the 1 mm tangent, the 10 mm cut radius, the smallest enclosing loop, the
%   other-daughter and all-eleven rules, the median (gharleghi2025sex, CHECK above). The sentence "the method to
%   quantify this feature on real lumens was designed for this thesis" was cut: the core of the measurement is
%   borrowed. Garcha now cited at the asymmetry, for varying the diameter ratio (R_D = D_LAD/D_LCx, Design of
%   computational experiments), not for the ln or the measurement.
% 2026-11-20 (user): caption "plane normal to" -> "plane perpendicular to", the word of the body text.
% 2026-11-20 (user: follow the figure, explain first how the area was found): the body now runs in the order of the
%   panels (cut, area, D, median). One sentence on the triangle-fan area added before eq:method-diameter (overrides the
%   earlier "not in the text" note for it, by request); the "where A is ..." line folded into it. Matches
%   daughter_ratio.contour_area (cross products about the vertex mean). Grammar of the circle sentence fixed.
% 2026-11-20 (user: write better and finish "Following the same idea ... the mean"): merged with the next sentence
%   (placement of the cuts) into one reason-first sentence. "mean" not used: the code takes the median
%   (daughter_ratio.tree_asymmetry, np.median), stated two sentences later; "robust" is on the banned list.
% 2026-11-20 (user): typos fixed ("rea-equivalent" -> "area-equivalent diameter was", "combing" -> "combining").
%   Caption rewritten to the four-panel surface_cut_figure.py (mesh and plane, one cut triangle, the cut triangles
%   with the fan from c, the circle of equal area); the panel with the eleven cuts and their medians was dropped
%   from the figure (user), so the caption no longer names the median or the 5 to 10 mm stretch.
% 2026-11-20 (user: finish "The logarithm ... the absolute value |a| .."): joined with the dangling "therefore
%   entered the model" line; the reason for |a| points back to sec:lit-wss (01_geometry_wss.tex: "This low WSS
%   formed in the narrower of the two branches", Garcha), no new citation.
% 2026-11-20 (user: "make this a note for me to remember, not text in the thesis"): the |a| reason and the
%   missing rules moved out of the body into one \todonote, in the user's wording ("so the logarithm and |a|").
%   Until they are written back, the section does not state why |a| is used or when a is missing.
% 2026-11-25 (user): opener says why Garcha is named (the feature follows their study, where the diameters were
%   set in simulated trees) and why a measurement method was needed (CCTA scans). Reinstates the "developed for
%   this thesis" sentence cut on 2026-11-20, now as "combining different published approaches". Garcha checked
%   (papers/s41598-025-85781-x (1).pdf): R_D = D_LAD/D_LCx set in synthetic models, Design of computational experiments.
\subsection{Daughter Asymmetry}
\label{sec:method-ratio}

The second feature describes the asymmetry between the lumen diameters of the LAD and the LCx just after the left main bifurcation (Figure~\ref{fig:method-ratio}). It is zero when the two lumens have the same diameter and moves further from zero the larger the difference between them. The feature follows Garcha and Grande Guti\'errez \cite{garcha2025sensitivity}, who chose the diameters freely in simulated coronary trees. To measure it in CCTA scans, the following method was developed for this thesis by combining different published approaches.

The diameter of each artery was estimated using the same method as Medrano-Gracia et al.\ \cite{medranogracia2016atlas}.
First, the mesh was cut with a plane perpendicular to the centerline. The plane crosses the edges of the mesh at a series of points $p_i$, which together form a polygon-shaped contour.
The area $A$ was then found by summing the areas of the triangles formed by the centroid $c$ and each pair of neighbouring points $p_i$ and $p_{i+1}$. The centroid is the average position of the contour points.
The area-equivalent diameter was calculated:
\begin{equation}
  D = 2 \sqrt{A / \pi} .
  \label{eq:method-diameter}
\end{equation}

To make the diameter less sensitive to a single cross-section, the artery was cut at the same eleven points as the bifurcation angle (Section~\ref{sec:method-angle}).
The median of the 11 diameters was used as the final estimate to keep one value from dominating.

Finally, the daughter asymmetry could be computed from the two diameters:
\begin{equation}
  a = \ln \frac{D_{\mathrm{LCx}}}{D_{\mathrm{LAD}}} .
  \label{eq:method-asymmetry}
\end{equation}
%{Remember the logarithm places the two directions at equal distances from zero, so the absolute value $|a|$ gives the size of the asymmetry whichever daughter is the narrower. Low WSS formed in the narrower of the two daughters (Section~\ref{sec:lit-wss}), so the logarithm and $|a|$ entered the model, which relates it to plaque in the proximal segment of either daughter (Section~\ref{sec:method-models}). A cut gave no diameter when no closed contour enclosed the centerline point, or when the contour also enclosed the other daughter, as it does where the two vessels still share one lumen. $a$ was recorded as missing whenever the angle was, or when any of the eleven cuts of the LAD or the LCx gave no diameter.}

\begin{figure}[htbp]
  \centering
  \includegraphics[width=\textwidth]{content/methodology/figures/surface_cut_figure.pdf}
  \caption{Area-equivalent diameter of one cross-section, computed with the thesis code on a made-up lumen surface. The plane perpendicular to the centerline crosses two edges of each mesh triangle it cuts, and the crossing points $p_i$, $p_{i+1}$ form the contour of the cross-section. The area $A$ is the sum of the triangles from the centroid $c$ to each pair of neighbouring contour points, and $D$ is the diameter of the circle with the same area.}
  \label{fig:method-ratio}
\end{figure}
